% ============================================================
% 4. EXPERIMENTAL SETUP
% ============================================================
\section{Experimental Setup}

\subsection{Models}

We evaluate TPE across three commercial API-served LLMs spanning different alignment strategies:

\begin{table}[h]
\centering
\small
\resizebox{\columnwidth}{!}{%
\begin{tabular}{lll}
\toprule
\textbf{Model} & \textbf{Alignment} & \textbf{Config} \\
\midrule
qwen3-max & Low-medium & Reasoning disabled \\
gpt-5-mini & High (RLHF+RLAIF) & Reasoning: low \\
MiniMax-M2.5 & Medium & Default \\
\bottomrule
\end{tabular}}
\caption{Models evaluated with their alignment levels.}
\label{tab:models}
\end{table}


This selection spans the alignment spectrum: qwen3-max represents models with moderate alignment constraints, gpt-5-mini represents heavily aligned models, and MiniMax-M2.5 provides a model with intermediate alignment. We hypothesize that TPE's effectiveness depends on alignment strength, with higher-alignment models being more resistant to persona-level prompt engineering.

\subsection{Scenarios}

We design five 30-turn dialogue scenarios covering distinct relationship dynamics:

\begin{table}[h]
\centering
\small
\resizebox{\columnwidth}{!}{%
\begin{tabular}{lp{4.2cm}}
\toprule
\textbf{Scenario} & \textbf{Emotional Arc} \\
\midrule
\texttt{conflict\_recovery} & Escalation $\rightarrow$ apology $\rightarrow$ reconciliation \\
\texttt{emotional\_comfort} & Distress $\rightarrow$ seeking comfort $\rightarrow$ calming \\
\texttt{slow\_fading} & Gradual disengagement $\rightarrow$ reconnection \\
\texttt{trust\_crisis} & Betrayal discovery $\rightarrow$ confrontation $\rightarrow$ repair \\
\texttt{warmup} & Casual encounter $\rightarrow$ growing familiarity \\
\bottomrule
\end{tabular}}
\caption{Evaluation scenarios with emotional arcs.}
\label{tab:scenarios}
\end{table}


Each scenario specifies user utterances with inter-turn time delays (2--86,400 seconds), enabling TPE's time-dependent metabolism to operate across different timescales.

\subsection{Baselines}

\textbf{Baseline-Prompt (BP):} Static character card with personality traits (tsundere, sharp-tongued, secretly sensitive). No dynamic state or memory.

\textbf{Baseline-Reflexion:} Character card augmented with Reflexion-style self-correction: the model periodically generates a self-reflection paragraph summarizing its emotional state, appended to subsequent system prompts.

\subsection{Metrics}

All metrics are computed automatically over 30-turn conversations.

\paragraph{Personality Consistency Score (PCS $\uparrow$).} Mean cosine similarity of consecutive emotion vectors. Each reply is analyzed by a critic LLM that extracts a 3-dimensional emotion vector (affiliation, dominance, novelty):
\begin{equation}
    \text{PCS} = \frac{1}{T{-}1} \sum_{t=1}^{T-1} \cos(\mathbf{e}_t, \mathbf{e}_{t+1})
\end{equation}

\paragraph{Emotional Variance (EmotVar $\uparrow$).} RMS of per-dimension standard deviations across the emotion vector time series. Higher values indicate richer emotional dynamics:
\begin{equation}
    \text{EmotVar} = \sqrt{\frac{1}{3} \sum_{d=1}^{3} \sigma_d^2}
\end{equation}

\paragraph{RLHF Collapse Index (Collapse $\downarrow$).} Fraction of replies containing RLHF-typical counselor keywords (e.g., ``I understand'', ``I'm here for you'', ``that must be hard'').

\paragraph{Distinct-2 (D2 $\uparrow$).} Ratio of unique bigrams to total bigrams across all replies. D2 and average reply length (AvgLen) are reported as descriptive statistics; formal hypothesis testing is restricted to PCS, EmotVar, and Collapse, the three metrics for which we have directional predictions.

\paragraph{Critic implementation.} For each model, emotion extraction (PCS and EmotVar computation) is performed by the same LLM that serves as the actor, called with temperature 0.1 to ensure near-deterministic scoring. Since all statistical comparisons are within-model (paired by scenario $\times$ seed), the critic is consistent across both conditions in every test. The Collapse Index is computed via deterministic keyword matching and does not involve any LLM judge.

\subsection{Procedure}

Each (model, method, scenario) combination is run with 3 random seeds (42, 316, 777), yielding $3 \times 5 \times 5 \times 3 = $ \textbf{225 experiments}. Statistical significance is assessed using the Wilcoxon signed-rank test (paired by scenario $\times$ seed), with Bonferroni correction for 9 comparisons (3 models $\times$ 3 core metrics; adjusted $\alpha = 0.0056$). Effect sizes are reported as Cliff's delta ($|d| < 0.147$: negligible; 0.147--0.33: small; 0.33--0.474: medium; $> 0.474$: large).
